\section{The Group}

There are two ways to understand the gauge group just obtained, and they make different physical predictions. Both are
laid out here, because the deciding observation has not yet been made.

On the first reading the gauge group is exactly the product of three pieces, read off three independent rungs of the
tower, and it is not the restriction of a larger single group, since the three factors come from three different division
algebras rather than one embedding. The consequence is sharp: the proton decay that a grand-unified gauge boson would
mediate is not predicted, because there are no such unifying bosons. A definite observation of proton decay in the classic
unified channels would sit awkwardly with this reading.

On the second reading the matter of one generation fills the sixteen-component spinor of a single larger group that holds
a full generation including a right-handed neutrino, and the Standard Model group sits inside it. The model's own
computation of the weak mixing angle uses this larger group's charge organization. On it,
grand-unified proton decay is generically allowed, with a lifetime fixed by the unification scale at about ten to the
thirty-six years, above the current bound and falsifiable.

The two readings agree on everything tested so far, the charges, the weak angle, the anomaly cancellation, and disagree
on one future observation: whether the proton decays through the grand-unified channels. That single measurement, or a
strong enough bound, decides between them. The representation theory that delivers the weak angle uses the larger
group's organization, while the dynamical content of the forces is the product. So the larger group organizes the
matter while the forces are the product, and whether it is also a dynamical gauge symmetry is the open fork.
